\section{Learning the launch footprint from a reference scene}
\label{sec:reference-calibration}
The footprint in the stationary experiment can itself be unknown.
A single calibrated reference obstacle, observed through the same
collision-bit interface, separates the footprint from the target
geometry. This is an additional calibration experiment, not a claim
that the original unreferenced data reveal an arbitrary unknown noise
support. Its preparations are included in the sample count.

\subsection{The reference experiment and its information}
Let $\mathcal K$ be a class of convex footprints with known uniform
$C^{6,\beta}$ support bounds in the laboratory frame, positive lower
and upper curvature-radius bounds, a known containing disk and
$\diam K\le\overline d_K<d_0$. The exact support function and the
placement of $K$ inside this disk are not supplied. The step satisfies
\eqref{eq:line-step}. Fix a known disk $R$ of positive radius at a
known laboratory location. In a separate calibration scene, $R$ is
the only solid in a known protective aperture. The same nominal
coordinate frame and the same footprint $K$ are used in both scenes.
Fresh independent launch draws have densities $j_R,j_T$, each stationary
within its scene and each obeying the same known lower bound
$b_0\dist(z,\partial K)^\gamma$. The two densities may differ.
Neither their values nor their means are measured. The calibration
apparatus must retain the footprint between the two scenes; this
assumption is not inferred by the controller from a finite record.

The two positive components in question have supports
\begin{equation}\label{eq:reference-addition}
 p_{P_R}=p_R+p_{-K},\qquad p_{P_C}=p_C+p_{-K}.
\end{equation}
All supports in this section are laboratory supports. A reference
with any known uniformly convex smooth support could replace the disk.
The known reference and its placement are metrological inputs. The
construction does not estimate their manufacturing cost.

\begin{theorem}[Joint footprint and table identification]
\label{thm:reference-exact}
Exact reciprocal pooled means in the reference and target apertures
determine $K$ and all protected original component bodies. On a
bounded periodic class they determine the full period group with exact
zero tests; a known nonperiod margin is needed only for uniform finite
decisions. In particular, neither $p_K$, its laboratory displacement,
nor either nuisance density is supplied to this inverse. The recovery
formulas are
\begin{equation}\label{eq:reference-subtraction}
 p_K(n)=p_{P_R}(-n)-p_R(-n),\qquad
 p_C(n)=p_{P_C}(n)-p_{P_R}(n)+p_R(n).
\end{equation}
For two such physical experiments, let $\Delta_R,\Delta_T$ be the
supremum differences of their reciprocal forcings on the respective
fixed apertures. For sufficiently small differences, the expanded
components have a bijection, and
\begin{align}
 \|p_{K_0}-p_{K_1}\|_\infty
       &\le C\Delta_R^{1/\alpha},\label{eq:reference-modulus-K}\\
 \max_C\|p_{C_0}-p_{C_1}\|_\infty
       &\le C(\Delta_T^{1/\alpha}+\Delta_R^{1/\alpha}),
                                           \label{eq:reference-modulus-C}\\
 \|p_{K_0}-p_{K_1}\|_{C^2}
    +\max_C\|p_{C_0}-p_{C_1}\|_{C^2}
       &\le C(\Delta_T^\zeta+\Delta_R^\zeta),\qquad
        \zeta=\frac{s-2}{s\alpha},\quad\alpha=\gamma+\frac32.
                                           \label{eq:reference-modulus-C2}
\end{align}
The densities and footprints may both differ between the two compared
experiments; the reference obstacle and the class bounds are common.
\end{theorem}
\begin{proof}
The stopped-walk construction and the support-positivity argument in
Lemma~\ref{lem:noise-support} recover the expanded components without
using the value of $p_K$. The reference scene contains just one such
component. Equations \eqref{eq:reference-addition} therefore give
\eqref{eq:reference-subtraction} and exact identification. Primitive
recognition is applied only after this subtraction.

The stopped comparison and the boundary-mass argument proving
Proposition~\ref{prop:jitter-modulus} do not require equality of
footprints before its final subtraction. Uniform footprint bounds
supply the same rolling radii, separated-component bounds and exit-time
constants. Hence the two reference expanded supports differ by at most
$C\Delta_R^{1/\alpha}$, and matched target expanded supports differ
by at most $C\Delta_T^{1/\alpha}$. Subtracting the common $p_R$
proves \eqref{eq:reference-modulus-K}; subtracting the reference
expanded support in each world proves \eqref{eq:reference-modulus-C}.
The signed-kernel interpolation estimate in
Proposition~\ref{prop:jitter-modulus}, applied on the uniform original
support classes, gives \eqref{eq:reference-modulus-C2}. The inequality
$(x+y)^{(s-2)/s}\le x^{(s-2)/s}+y^{(s-2)/s}$ separates the two
error sources. No density deconvolution, equality of densities, or
unknown footprint evaluation has entered the proof.
\end{proof}

\begin{theorem}[Finite self-calibration with all attempts charged]
\label{thm:reference-finite}
On $\mathcal B_\eta\times\mathcal K$, with the reference experiment
just specified, a finite pooled-bit controller reconstructs $K$ and
the original periodic table to $C^2$ accuracy $C\nu$, and recovers
its primitive discrete data, with probability at least $1-\delta$.
It does not use an oracle for $K,j_R$ or $j_T$. Its total number of
attempts, including reference-scene preparations, is at most
\begin{equation}\label{eq:reference-cost}
 C\nu^{-((\gamma+2)s+1)/(s-2)}
                   \log(C/\nu)\log(C/(\nu\delta)).
\end{equation}
The apparatus makes one fixed scene change and does not reduce the
launch footprint with $\nu$. The nominal-coordinate accuracy,
line-batch counts, spatial-grid counts and numerical qualifications
are those of Theorem~\ref{thm:line-reconstruction}. Known positive
geometric and nonperiod margins, exact compass flights, and uniform
footprint bounds remain essential assumptions.
\end{theorem}
\begin{proof}
Apply the coarse phase and the one-sided support procedure to the
reference expanded body $P_R$ and to all complete target expanded
bodies $P_C$. These procedures require only bounds on their support
norms, radii, location and separation. Bounds for $R$ and $\mathcal K$
supply them in the reference scene; bounds for the target and
$\mathcal K$ supply them in the other scene. No prescribed support
function of $K$ is used in the coarse search or in the line response.
For $h\asymp\nu^{1/(s-2)}$ and $e=c h^s$, obtain smooth estimates
of these expanded supports with errors $C h^s$ in supremum norm and
$C h^{s-2}$ in $C^2$. There is one additional expanded body, not an
accuracy-dependent number of reference parameters.

Define
\[
 \widehat p_K(n)=\widehat p_{P_R}(-n)-p_R(-n),\qquad
 \widehat p_C(n)=\widehat p_{P_C}(n)-\widehat p_{P_R}(n)+p_R(n).
\]
The two error bounds add with a fixed factor. The positive lower
curvature radii of $K$ and $C$ ensure that these are supports of actual
convex bodies for sufficiently small $\nu$, after the same finite
smooth-output construction used earlier. Their Hausdorff errors remain
$O(h^s)$. Apply Lemma~\ref{lem:locking} to the reconstructed original
components; it does not require the footprint or reference to have any
period. Its known thresholds give the primitive discrete decisions,
and real basis vectors and free area have the stated geometric accuracy.

Allocate failure probability over both coarse phases and all line
batches before the experiment. The target component count in the
protected patch is bounded by the priors. Adding one reference body
therefore changes only the constant in \eqref{eq:line-total}.
Conditional fresh sampling and the union bound prove
\eqref{eq:reference-cost}. A direction-resolved response, observed
actual launch position, and density-estimation stage are not required.
\end{proof}

\begin{remark}[Reference error is not erased by repetition]
\label{rem:reference-error}
If the reference support used in the subtraction differs from the
actual one by $\varepsilon_j$ in $C^j$, then the recovered target
supports have error at most
$C(h^{s-j}+e h^{-j})+\varepsilon_j$ for $j=0,2$.
This follows directly by subtracting \eqref{eq:reference-subtraction};
the reference error enters once with coefficient one. Curvature and
period conclusions require these total errors below their corresponding
margins. An unknown reference translation therefore remains an absolute
position uncertainty, not a vanishing noise term. The reference theorem
trades complete prior knowledge of the footprint for a calibrated
physical reference, while measuring the footprint with collision bits.
It does not remove every form of calibration or prove blind recovery
from the original target scene alone.
\end{remark}
